% ch3.tex -- Module 7 ESSENTIALS, Chapter 3: Light as a particle (black bodies and
% the photoelectric effect). Knowledge points only: no questions, no worked examples.
% Figures: figures/ch3_bbcurves.pdf and figures/ch3_uvcat.pdf (copied from
% ../mod7/figures, with their scripts); the other diagrams are TikZ copied from the
% full chapter.
\chapter{Light as a Particle: Black Bodies and the Photoelectric Effect}\label{ch3:chap}
\covers{M7.IQ3.01, M7.IQ3.02, M7.IQ3.03}

\begin{summarybox}{In this chapter}
\begin{itemize}
  \item Black body curves and Wien's law.
  \item The ultraviolet catastrophe, and how Planck's quanta, $E = hf$, solved it.
  \item The photoelectric results that the wave model of light cannot explain.
  \item Einstein's photons, $K_\mathrm{max} = hf - \phi$, and photoelectric graphs.
\end{itemize}
\end{summarybox}

%% ==================================================================
\section{Black body radiation and Planck's quanta}\label{ch3:sec:blackbody}
\covers{M7.IQ3.01}

\begin{bigidea}
A black body's spectrum depends only on its temperature. In 1900 Planck explained its
shape: energy is given out in packets, $E = hf$.
\end{bigidea}

\begin{keywords}
  \kw{Black body}{an ideal object that absorbs all the radiation that hits it. It is also a perfect emitter.}
  \kw{Peak wavelength ($\lambda_\mathrm{max}$)}{the wavelength where a black body curve is highest.}
  \kw{Ultraviolet catastrophe}{the wrong prediction, from physics before 1900, of infinite intensity at short wavelengths.}
  \kw{Quantum (plural: quanta)}{one packet of energy, of size $E = hf$.}
\end{keywords}

\subsection{The black body curve and Wien's law}

\begin{itemize}
  \item All objects give off \term{thermal radiation} because of their temperature:
        their vibrating charges give off EM waves.
  \item A hollow box with a tiny hole is almost a black body. So are stars.
  \item Lava (about \qtyrange{1000}{1500}{K}) peaks in the infrared and glows dull red.
        The Sun (\qty{5800}{K}) looks white. Hotter stars look blue-white.
  \item Wien (1893): $\lambda_\mathrm{max}$ is \emph{inversely proportional} to $T$. If
        $T$ doubles, $\lambda_\mathrm{max}$ halves.
  \item \term{Intensity}: radiation energy each second on each square metre
        (\cref{ch3:fig:bbcurves}).
\end{itemize}

\begin{sidebyside}[0.57]
  \centering
  \includegraphics{ch3_bbcurves.pdf}
  \captionof{figure}{Black body curves at three temperatures.}\label{ch3:fig:bbcurves}
\nextside
  \begin{itemize}
    \item One peak, at $\lambda_\mathrm{max}$. Steep rise on the short-wavelength side,
          long tail on the long side.
    \item A hotter body is higher at \emph{every} wavelength. The curves never cross.
    \item A hotter body peaks at a shorter wavelength.
    \item The curve depends only on temperature, not on the material.
  \end{itemize}
\end{sidebyside}

\begin{formulacard}{Wien's displacement law}
  \cardformula{\lambda_\mathrm{max} = \frac{b}{T}}
  \begin{symbols}
    \sym{$\lambda_\mathrm{max}$}{peak wavelength of the black body curve}{\si{m}}
    \sym{$b$}{Wien's displacement constant, \qty{2.898e-3}{m.K}}{\si{m.K}}
    \sym{$T$}{absolute temperature of the object}{\si{K}}
  \end{symbols}
  \whentouse{black bodies or near black bodies: stars, lava, hot metal, people.}
  \trap{$T$ in kelvin (add 273 to \si{\degreeCelsius}). $\lambda_\mathrm{max}$ comes out in
    metres ($\qty{1}{nm} = \qty{1e-9}{m}$).}
\end{formulacard}

\subsection{The ultraviolet catastrophe}

\begin{sidebyside}[0.56]
  \centering
  \includegraphics{ch3_uvcat.pdf}
  \captionof{figure}{Classical prediction and measured curve.}
  \label{ch3:fig:uvcat}
\nextside
  \begin{itemize}
    \item \term{Classical physics} (before 1900): light is a wave and energy can have any
          value.
    \item Its prediction (the \term{Rayleigh--Jeans law}) fits at long wavelengths
          (\cref{ch3:fig:uvcat}).
    \item At short wavelengths it rises without limit, so the total energy would be
          infinite. Real curves fall to zero.
  \end{itemize}
\end{sidebyside}

Why it failed: a hot box holds more and more standing-wave patterns at shorter
wavelengths. Classical physics gives each pattern the same share of energy, so the
total is infinite.

\subsection{Planck's idea: energy comes in packets}

\begin{itemize}
  \item Planck (1900): the vibrating charges in the walls (\term{oscillators}) have
        \term{quantised} energy. It comes in steps: only $0, hf, 2hf, 3hf, \ldots$
  \item So they give out or take in energy only in quanta of $hf$. $h$ is
        \term{Planck's constant}, \qty{6.626e-34}{J.s}.
  \item Low $f$ (long wavelength): the steps are tiny and act like a ramp. Planck agrees
        with classical physics there.
  \item High $f$ (short wavelength): one step $hf$ is bigger than the heat energy an
        oscillator usually has. Few quanta are given out, so the curve falls to zero.
  \item Planck's formula fits the measured curves at every wavelength. He saw quanta as
        a maths trick, but they began the quantum model of light.
\end{itemize}

\begin{formulacard}{Energy of one quantum (photon)}
  \cardformula{E = hf}
  \begin{symbols}
    \sym{$E$}{energy of one quantum (one photon)}{\si{J}}
    \sym{$h$}{Planck's constant, \qty{6.626e-34}{J.s} (data sheet)}{\si{J.s}}
    \sym{$f$}{frequency of the radiation}{\si{Hz}}
  \end{symbols}
  \whentouse{the energy of one quantum or one photon of any EM radiation.}
  \derive{$E = \dfrac{hc}{\lambda}$ is not on the sheet. Use $v = f\lambda$ with $v = c$
    (speed of light): $f = \dfrac{c}{\lambda}$, so $E = \dfrac{hc}{\lambda}$.}
  \trap{$E$ comes out in joules. The \term{electron volt}: $\qty{1}{eV} =
    \qty{1.602e-19}{J}$ (data sheet). J $\to$ eV: divide by \num{1.602e-19}.}
\end{formulacard}

%% ==================================================================
\section{The photoelectric effect: evidence against the wave model}\label{ch3:sec:photoelectric}
\covers{M7.IQ3.02}

\begin{bigidea}
Light can knock electrons out of a metal. Experiments showed that three predictions
of the wave model of light are wrong.
\end{bigidea}

\begin{keywords}
  \kw{Photoelectric effect}{electrons are released from a metal surface when light shines on it.}
  \kw{Threshold frequency ($f_0$)}{the lowest frequency of light that releases electrons from a metal.}
  \kw{Stopping voltage ($V_0$)}{the smallest reverse voltage that stops every photoelectron, so the current is zero.}
\end{keywords}

\begin{itemize}
  \item Hertz first saw the effect (1887): ultraviolet light helped sparks jump a
        \term{spark gap} (a small air gap between two metal balls). Lenard measured it with
        a vacuum tube (1902).
\end{itemize}

\begin{sidebyside}[0.62]
  \centering
  \begin{tikzpicture}[x=0.9cm, y=0.82cm]
    \draw[hscink, line width=0.9pt, fill=hscbrand!5, rounded corners=12pt]
      (0,0) rectangle (6,2.6);
    \draw[hscvel!70, line width=2.6pt] (1.15,2.6) -- (2.75,2.6);
    \node[diagramlabel, anchor=south east, align=right, font=\footnotesize] at (1.1,2.65)
      {quartz window\\(lets UV through)};
    \draw[hscink, line width=2.4pt] (0.8,0.45) -- (0.8,2.15);
    \draw[hscink, line width=2.4pt] (5.2,0.6) -- (5.2,2.0);
    \foreach \ye/\xs in {1.9/2.42, 1.45/2.78, 1.0/3.14} {
      \draw[->, hscray, line width=1pt, decorate,
        decoration={snake, amplitude=1.3pt, segment length=5pt, post length=2.2mm}]
        (\xs,3.8) -- (0.92,\ye);
    }
    \node[diagramlabel, anchor=west, font=\footnotesize] at (3.2,3.75) {light: frequency $f$};
    \foreach \p in {(1.9,1.65), (3.0,1.0), (4.0,1.75)} {
      \fill[hscvel] \p circle (2.6pt);
      \draw[velocity, line width=0.9pt] \p ++(0.12,0) -- ++(0.55,0);
    }
    \node[text=hscvel, font=\footnotesize] at (3.0,0.6) {photoelectrons};
    \node[diagramlabel, anchor=north, font=\footnotesize] at (3.6,2.45) {vacuum};
    \draw[hscmuted, line width=0.5pt] (0.75,1.6) -- (-0.2,1.6)
      node[diagramlabel, anchor=east, align=right, font=\footnotesize] {metal\\cathode};
    \draw[hscmuted, line width=0.5pt] (5.25,1.6) -- (6.2,1.6)
      node[diagramlabel, anchor=west, align=left, font=\footnotesize] {anode};
    % circuit
    \draw[hscink, line width=0.8pt] (0.8,0.45) -- (0.8,-1.0) -- (2.2,-1.0);
    \draw[hscink, line width=0.8pt] (2.4,-1.0) -- (3.67,-1.0);
    \draw[hscink, line width=0.8pt] (4.33,-1.0) -- (5.2,-1.0) -- (5.2,0.6);
    \draw[hscink, line width=0.9pt] (2.2,-1.35) -- (2.2,-0.65);
    \draw[hscink, line width=2.2pt] (2.4,-1.18) -- (2.4,-0.82);
    \draw[->, hscink, line width=0.7pt] (1.9,-1.4) -- (2.7,-0.55);
    \node[font=\scriptsize] at (2.02,-0.5) {$+$};
    \node[font=\scriptsize] at (2.6,-1.3) {$-$};
    \draw[hscink, line width=0.8pt] (1.6,-1.0) -- (1.6,-1.8) -- (2.0,-1.8);
    \draw[hscink, line width=0.8pt] (2.6,-1.8) -- (3.0,-1.8) -- (3.0,-1.0);
    \draw[hscink, line width=0.8pt, fill=white] (2.3,-1.8) circle (0.3) node[font=\small] {V};
    \node[diagramlabel, anchor=north, font=\footnotesize, align=center] at (2.3,-2.15)
      {reversible supply and voltmeter};
    \draw[hscink, line width=0.8pt, fill=white] (4.0,-1.0) circle (0.34)
      node[font=\small] {\si{\micro A}};
    \node[diagramlabel, anchor=south, font=\footnotesize] at (4.0,-0.64) {microammeter};
  \end{tikzpicture}
  \captionof{figure}{Photoelectric apparatus. Here the anode is negative, as used to
    find $V_0$.}\label{ch3:fig:apparatus}
\nextside
  \begin{itemize}
    \item Light releases \term{photoelectrons} from the \term{cathode}. They cross to
          the \term{anode}.
    \item A \term{microammeter} (a meter for tiny currents) measures the
          \term{photocurrent}. Bigger current $=$ more electrons each second, not
          faster ones.
    \item Vacuum: electrons do not hit air molecules. Quartz lets UV in (glass
          absorbs it).
  \end{itemize}
\end{sidebyside}

\begin{itemize}
  \item \textbf{Electroscope demonstration:} an \term{electroscope} has a thin gold leaf
        that stands out when it is charged. Put a negatively charged zinc plate on it. UV
        light makes the leaf fall: electrons leave the zinc. Bright visible light does not.
\end{itemize}

\subsection{Stopping voltage}

\begin{itemize}
  \item Make the anode negative (\cref{ch3:fig:apparatus}). It repels the electrons.
        At $V_0$ even the fastest ones stop just before the anode, so $I = 0$.
  \item Electrons leave with a range of speeds. The fastest have the
        \term{maximum kinetic energy}, $K_\mathrm{max}$.
  \item $K_\mathrm{max} = qV_0$ \hscbadgecombine{$W = qV$}\ The reverse voltage does work $W = qV_0$ against the fastest electron and takes
        all its kinetic energy. $q$ is the \emph{size} of the electron's charge,
        \qty{1.602e-19}{C}.
  \item $V_0$ is a voltage (V); $K_\mathrm{max}$ is an energy (J or eV). In eV,
        $K_\mathrm{max}$ has the same number as $V_0$ in volts.
  \item \textbf{Speed of the fastest electrons:} put $K_\mathrm{max}$, in joules, into
        $K = \frac{1}{2}mv^2$ (on the sheet), with $m$ = the electron's mass.
\end{itemize}

\subsection{What the wave model predicted}

In the wave model, light energy depends on brightness (amplitude), not frequency. It
arrives smoothly, spread over the surface (\cref{ch3:tab:wave}).

\begin{table}
  \caption{Wave-model predictions and photoelectric observations.}\label{ch3:tab:wave}
  \centering
  \small
  \begin{tabular}{@{}L{52mm}L{64mm}L{20mm}@{}}
    \toprule
    \colhead{The wave model predicted} & \colhead{Experiments showed} & \colhead{Wave model}\\
    \midrule
    Brighter light gives more electrons each second (bigger current).
      & Brighter light gives a bigger current (when $f \ge f_0$).
      & \cmark\ agrees\\
    Any frequency works if the light is bright enough.
      & Threshold frequency $f_0$: below it, no electrons, however bright.
      & \xmark\ fails\\
    Dim light: a delay while electrons build up energy.
      & No delay: electrons leave at once (in less than \qty{1e-9}{s}).
      & \xmark\ fails\\
    Brighter light makes faster electrons (bigger $V_0$).
      & $V_0$ does not depend on brightness. It rises in a straight line with $f$.
      & \xmark\ fails\\
    \bottomrule
  \end{tabular}
\end{table}

So light energy cannot arrive spread out like a wave. It must arrive in concentrated
packets of $hf$.

%% ==================================================================
\section{Einstein's photon model and photoelectric graphs}\label{ch3:sec:photon}
\covers{M7.IQ3.03, M7.IQ3.02}

\begin{bigidea}
In 1905 Einstein said that light itself is a stream of photons, each with energy
$hf$. One photon gives all its energy to one electron, so conservation of energy gives
$K_\mathrm{max} = hf - \phi$.
\end{bigidea}

\begin{keywords}
  \kw{Photon}{one packet (quantum) of light, with energy $E = hf$.}
  \kw{Work function ($\phi$)}{the least energy to remove an electron from a metal's surface.}
  \kw{Saturation current}{the largest photocurrent: all released electrons reach the anode.}
  \kw{Wave--particle duality}{light acts as a wave in some experiments, as particles in others.}
\end{keywords}

\begin{itemize}
  \item Brighter light of one colour has \emph{more photons} each second. Each photon
        still has the same energy $hf$.
  \item \textbf{One photon interacts with one electron}, and gives it all its energy at
        once. Energy from two photons never adds up.
  \item Each metal has its own $\phi$, for example caesium \qty{2.1}{eV}, zinc
        \qty{4.3}{eV}.
\end{itemize}

\begin{formulacard}{Einstein's photoelectric equation}
  \cardformula{K_\mathrm{max} = hf - \phi}
  \begin{symbols}
    \sym{$K_\mathrm{max}$}{maximum kinetic energy of the photoelectrons}{\si{J} (or eV)}
    \sym{$h$}{Planck's constant, \qty{6.626e-34}{J.s}}{\si{J.s}}
    \sym{$f$}{frequency of the light}{\si{Hz}}
    \sym{$\phi$}{work function of the metal}{\si{J} (or eV)}
  \end{symbols}
  \whentouse{light falls on a metal and $hf \ge \phi$. If $hf - \phi$ is negative, no
    electrons are released.}
  % layout (final page check): break the card here, so that the page after it does not
  % hold only the two-line Trap
  \tcbbreak
  \derive{conservation of energy for one photon and the \term{least-bound electron}
    (the easiest to remove): energy in $=$ energy to escape $+$ energy left over. So
    $hf = \phi + K_\mathrm{max}$, and $K_\mathrm{max} = hf - \phi$. Deeper electrons
    need more than $\phi$, so they leave slower.}
  \trap{use one unit for every term: all joules or all eV. $\phi$ is lower-case phi;
    capital $\Phi$ on the sheet is magnetic flux.}
\end{formulacard}

The photon model explains every result in \cref{ch3:tab:wave}:
\begin{itemize}
  \item \textbf{Bigger current:} more photons each second free more electrons each second.
  \item \textbf{Threshold:} one photon must give at least $\phi$, so $hf \ge \phi$.
  \item \textbf{No delay:} each photon delivers its whole energy in one event.
  \item \textbf{$K_\mathrm{max}$ depends on $f$ (and the metal), not on brightness:}
        brightness changes the number of photons, not $hf$.
\end{itemize}

\begin{itemize}
  \item \textbf{At the threshold} $f_0$ the fastest electron has nothing left over. Put
        $K_\mathrm{max} = 0$ in $K_\mathrm{max} = hf - \phi$:
        \[\phi = hf_0 \qquad \lambda_0 = \frac{c}{f_0} = \frac{hc}{\phi}\]
        {\centering\hscbadgenotonsheet\par}
  \item The \term{threshold wavelength} $\lambda_0$ is the longest wavelength that
        releases electrons. Electrons leave only if $f \ge f_0$ ($\lambda \le
        \lambda_0$), however bright the light.
  \item \textbf{Photons each second} \hscbadgenotonsheet\ from a source of power
        $P$: $\dfrac{N}{t} = \dfrac{P}{hf}$ (energy each second $\div$ energy of one
        photon). At the same power, red light gives more photons
        than blue.
\end{itemize}

\subsection{Planck, Einstein and a new model of light}

\begin{itemize}
  \item \textbf{Planck (1900)} quantised the oscillators in the walls. He saw quanta as a
        maths trick. They explained the black body curve.
  \item \textbf{Einstein (1905)} quantised light itself: photons are real. They explained
        the photoelectric effect. Both used $E = hf$.
  \item Light still shows interference and diffraction, but it is given out and absorbed
        as photons: wave--particle duality. (Photons are \emph{not} Newton's
        corpuscles.)
\end{itemize}

\subsection{Photoelectric graphs}

$K_\mathrm{max} = hf - \phi$ gives a straight-line graph of $K_\mathrm{max}$ against $f$,
with gradient $h$ and $K_\mathrm{max}$-intercept $-\phi$ (\cref{ch3:fig:kf}).

\begin{figure}
  \centering
  \begin{tikzpicture}
    \begin{axis}[hsc, width=0.74\linewidth, height=0.44\linewidth,
        axis lines=middle, xmin=0, xmax=14.5, ymin=-5, ymax=4.2,
        xtick={2,4,...,14}, xticklabels={}, ytick={-4,-2,2,4},
        xlabel={$f$ ($\times10^{14}$ Hz)}, ylabel={$K_\mathrm{max}$ (eV)},
        x label style={anchor=west}, y label style={anchor=south, rotate=-90},
        every axis x label/.style={at={(axis cs:14.5,0)}, anchor=west},
        every axis y label/.style={at={(axis cs:0,4.2)}, anchor=south},
        yticklabel style={font=\footnotesize}]
      \addplot[hscseriesA, line width=1.4pt, domain=5.077:14, samples=2] {0.4136*x - 2.10};
      \addplot[hscseriesA, line width=1pt, densely dotted, domain=0:5.077, samples=2] {0.4136*x - 2.10};
      \addplot[hscseriesB, line width=1.4pt, domain=10.40:14, samples=2] {0.4136*x - 4.30};
      \addplot[hscseriesB, line width=1pt, densely dotted, domain=0:10.40, samples=2] {0.4136*x - 4.30};
      \fill[hscseriesA] (axis cs:5.077,0) circle (2.2pt);
      \fill[hscseriesB] (axis cs:10.40,0) circle (2.2pt);
      \fill[hscseriesA] (axis cs:0,-2.10) circle (2.2pt);
      \fill[hscseriesB] (axis cs:0,-4.30) circle (2.2pt);
      \node[anchor=south east, font=\footnotesize, text=hscseriesA] at (axis cs:5.0,0.1) {$f_0$};
      \node[anchor=south east, font=\footnotesize, text=hscseriesB] at (axis cs:10.3,0.1) {$f_0$};
      \node[anchor=north west, font=\footnotesize, text=hscseriesA, fill=white, inner sep=1pt] at (axis cs:0.2,-2.2) {$-\phi = \qty{-2.10}{eV}$};
      \node[anchor=north west, font=\footnotesize, text=hscseriesB] at (axis cs:0.2,-4.4) {$-\phi = \qty{-4.30}{eV}$};
      \node[anchor=west, font=\small] at (axis cs:14.1,3.69) {caesium};
      \node[anchor=west, font=\small] at (axis cs:14.1,1.49) {zinc};
      \node[anchor=west, font=\footnotesize, align=left] at (axis cs:6.4,-3.4)
        {same gradient $h$: parallel lines.\\dotted: extrapolated (no electrons)};
      % f tick labels drawn after the lines, on white, so no line crosses a number
      \pgfplotsinvokeforeach{2,4,6,8,10,12,14}{
        \node[anchor=north, font=\footnotesize, fill=white, inner sep=1pt, yshift=-2pt]
          at (axis cs:#1,0) {#1};}
    \end{axis}
  \end{tikzpicture}
  \caption{$K_\mathrm{max}$ against frequency for caesium and zinc.}\label{ch3:fig:kf}
\end{figure}

\begin{itemize}
  \item \textbf{Gradient} $= h$ for every metal, so the lines are parallel
        (\qty{4.136e-15}{eV.s} with $K_\mathrm{max}$ in eV).
  \item \textbf{$f$-intercept} $= f_0$. \textbf{$K_\mathrm{max}$-intercept} $= -\phi$:
        \term{extrapolate} (extend the line, dotted) to $f = 0$, then change the sign.
  \item Experiments measure $V_0$. Put $K_\mathrm{max} = qV_0$ and divide by $q$:
        \[V_0 = \frac{h}{q}f - \frac{\phi}{q}\]
        {\centering\hscbadgenotonsheet\par}
        Gradient $\dfrac{h}{q}$ (so $h = q \times \text{gradient}$), $f$-intercept
        $f_0$, $V_0$-intercept $-\dfrac{\phi}{q}$.
\end{itemize}

\begin{sidebyside}[0.47]
  \centering
  \begin{tikzpicture}
    \begin{axis}[hsc, width=\linewidth, height=0.78\linewidth,
        axis lines=middle, xmin=-3, xmax=1.8, ymin=0, ymax=3.4,
        xtick={-2.4,-1.2}, xticklabels={$-V_{0}$ (A), $-V_{0}$ (B)}, ytick=\empty,
        xlabel={$V$}, ylabel={$I$ (photocurrent)},
        every axis x label/.style={at={(axis cs:1.8,0)}, anchor=west},
        every axis y label/.style={at={(axis cs:0,3.4)}, anchor=south},
        xticklabel style={font=\footnotesize}]
      \addplot[hscseriesA, line width=1.3pt, domain=-2.4:0, samples=40] {2.2*sin(90*(x+2.4)/2.4)^2};
      \addplot[hscseriesA, line width=1.3pt, domain=0:1.6, samples=2] {2.2};
      \addplot[hscseriesB, line width=1.3pt, dashed, domain=-1.2:0, samples=40] {1.1*sin(90*(x+1.2)/1.2)^2};
      \addplot[hscseriesB, line width=1.3pt, dashed, domain=0:1.6, samples=2] {1.1};
      \node[anchor=south, font=\footnotesize] at (axis cs:1.62,2.2) {A};
      \node[anchor=south, font=\footnotesize] at (axis cs:1.62,1.1) {B};
      \node[anchor=north, font=\scriptsize, align=center] at (axis cs:0.85,0.9)
        {saturation\\current (B)};
      \draw[->, hscink, line width=0.5pt] (axis cs:0.85,0.9) -- (axis cs:0.85,1.05);
      \node[anchor=south, font=\scriptsize, align=center, text=hscmuted]
        at (axis cs:-1.5,2.55) {anode negative\\(pushes electrons back)};
      \node[anchor=south, font=\scriptsize, align=center, text=hscmuted]
        at (axis cs:0.85,2.55) {anode\\positive};
    \end{axis}
  \end{tikzpicture}
  \captionof{figure}{Photocurrent against tube voltage (idealised).}
  \label{ch3:fig:iv}
\nextside
  \begin{itemize}
    \item Anode positive: all electrons are collected. The flat part (saturation
          current) is set by the number of photons each second.
    \item Beam A (\cref{ch3:fig:iv}) has more photons each second and a higher $f$
          (bigger $V_0$) than beam B.
    \item Brighter light, same $f$: only the flat part rises.
    \item Bigger $\phi$, same light: smaller $V_0$.
  \end{itemize}
\end{sidebyside}

%% ==================================================================
\section{Key experiments}\label{ch3:sec:experiments}
\covers{M7.IQ3.01, M7.IQ3.02, M7.IQ3.03}

\begin{experiment}{Measuring black body radiation (1899--1900)}
\aim{To measure black body curves at several temperatures (Lummer and Pringsheim).}
\method{Radiation leaves a small hole in a hot oven (the black body). A thermocouple (a
  thermometer) measures the oven's temperature. The radiation passes through a slit and a
  prism (\cref{ch3:fig:bbapp}). A \term{bolometer} (a detector warmed by the
  radiation) is moved along the spectrum. This is repeated at several temperatures.
  \begin{figure}
    \centering
    \begin{tikzpicture}[x=1cm, y=1cm]
      \fill[hscmuted!40] (0,0) rectangle (2.3,1.7);
      \fill[hscink!85] (0.35,0.35) rectangle (1.95,1.35);
      \fill[white] (1.94,0.75) rectangle (2.31,0.95);
      \draw[hscforce, line width=0.9pt] (0.9,1.9) -- (0.9,1.0);
      \node[diagramlabel, anchor=south, font=\footnotesize] at (0.9,1.9) {thermocouple};
      \node[diagramlabel, anchor=north, font=\footnotesize] at (1.15,-0.05)
        {hot oven (black body)};
      \draw[hscink, line width=2pt] (3.0,0.35) -- (3.0,0.78) (3.0,0.92) -- (3.0,1.35);
      \node[diagramlabel, anchor=north, font=\footnotesize] at (3.0,0.3) {slit};
      \draw[ray] (2.3,0.85) -- (4.1,0.85);
      \draw[hscink, line width=0.8pt, fill=hscbrand!10] (3.8,0.25) -- (5.0,0.25) -- (4.4,1.45) -- cycle;
      \draw[hscray, line width=1pt] (4.1,0.85) -- (4.73,0.79);
      \node[diagramlabel, anchor=north, font=\footnotesize] at (4.4,0.2) {prism};
      \draw[hscseriesB, line width=1pt] (4.73,0.80) -- (7.4,0.55);
      \draw[hscseriesC, line width=1pt] (4.73,0.79) -- (7.4,0.15);
      \draw[hscvel, line width=1pt] (4.73,0.78) -- (7.4,-0.25);
      \node[diagramlabel, anchor=west, font=\footnotesize] at (7.45,0.55) {longer $\lambda$ (bent least)};
      \node[diagramlabel, anchor=west, font=\footnotesize] at (7.45,-0.25) {shorter $\lambda$ (bent most)};
      \draw[hscink, line width=0.8pt, fill=white] (6.35,0.21) rectangle (6.65,0.51);
      \draw[hscmuted, line width=0.5pt] (6.5,0.51) -- (6.2,1.9)
        node[diagramlabel, anchor=south, font=\footnotesize] {bolometer (moved along the spectrum)};
    \end{tikzpicture}
    \caption{Measuring a black body curve (simplified).}\label{ch3:fig:bbapp}
  \end{figure}}
\results{Curves like \cref{ch3:fig:bbcurves}; $\lambda_\mathrm{max}T \approx
  \qty{2.9e-3}{m.K}$ at every temperature.}
\conclusion{Wien's law. Classical theory fitted only the long-wavelength side.}
\significance{Planck's formula fitted the whole curve. This led him to quanta (1900).}
\end{experiment}

\begin{experiment}{The photoelectric effect: measuring \texorpdfstring{$h$}{h} and \texorpdfstring{$\phi$}{phi}}
\aim{To find how $V_0$ depends on frequency and brightness (Lenard 1902, Millikan
  1914--1916, or a school photocell).}
\method{In the circuit of \cref{ch3:fig:apparatus}, raise the reverse voltage until
  the current is just zero ($V_0$). Repeat for several frequencies and brightnesses.}
\results{$V_0$ rises in a straight line with $f$. No current below $f_0$. Brightness
  changes the current, not $V_0$.}
\conclusion{$qV_0 = hf - \phi$, so $h = q \times \text{gradient}$; the intercepts give $f_0$ and
  $\phi$. This supports the photon model.}
\significance{Millikan (1914--1916) expected to prove Einstein wrong. Instead he
  confirmed the law and found $h$ within about 1\,\% of today's value.}
\end{experiment}

\begin{summarybox}{Remember}
\begin{itemize}
  \item A hotter black body is higher at every wavelength and peaks at a shorter
        wavelength: $\lambda_\mathrm{max} = \dfrac{b}{T}$, with $T$ in kelvin.
  \item Classical physics predicted the ultraviolet catastrophe. Planck: energy comes in
        quanta $E = hf$, and large high-$f$ quanta are rarely given out.
  \item Planck quantised the oscillators; Einstein quantised light itself (photons).
  \item The wave model fails: threshold frequency, no delay, $K_\mathrm{max}$ does not
        depend on brightness.
  \item One photon $\to$ one electron: $K_\mathrm{max} = hf - \phi = qV_0$. Brighter
        light gives more electrons, not faster ones.
  \item $K_\mathrm{max}$ against $f$: gradient $h$, $f$-intercept $f_0$, intercept
        $-\phi$.
\end{itemize}
\end{summarybox}
